% ch03 —— §3 超凯勒度量的扭量构造（原文 wcf-paper.tex 1165–1388 行）
\section{超凯勒度量的扭量构造}
\label{sec:twistor}

本节先回顾超凯勒几何的一些一般事实，然后说明我们对 $g$ 的描述背后的基本想法。

\subsection{来自超凯勒流形的全纯数据}
\label{sec:hk-general}

我们首先回顾附着在任一超凯勒流形上的一些全纯数据。按定义，超凯勒流形 $(\CM,g)$ 关于复结构的三元组 $\vec{J}$ 是凯勒的，这些复结构满足关系
\begin{equation} \label{eq:quaternion-complex-structures}
J_1 J_2 = J_3, \quad J_2 J_3 = J_1, \quad J_3 J_1 = J_2, \quad
J_\alpha^2 = -1.
\end{equation}
记 $\omega_\alpha$ 为对应的三个凯勒形式。

事实上，任一超凯勒流形 $(\CM,g)$ 关于更一般的复结构也是凯勒的，即关于 $a^\alpha J_\alpha$ 是凯勒的，其中 $\sum_{\alpha=1}^3 a_\alpha^2 = 1$，对应的凯勒形式为 $a^\alpha \omega_\alpha$。于是我们手上有整整一个 $S^2$ 的复结构。扭量进路的关键洞见之一在于：把这个 $S^2$ 看作黎曼球面、并用一个复参数 $\zeta$ 来标记，是有益的。因此我们把一般的复结构与对应的凯勒形式写成
\begin{align} \label{eq:Jzeta}
J^\pz &= \frac{i(- \zeta + \bar\zeta)J_1 - (\zeta + \bar\zeta) J_2 + (1 - \abs{\zeta}^2) J_3}{1 + \abs{ \zeta}^2},\\
\omega^\pz &= \frac{i(- \zeta + \bar\zeta)\omega_1 -  (\zeta + \bar\zeta) \omega_2 + (1 - \abs{\zeta}^2)\omega _3}{1 + \abs{\zeta}^2}.
\end{align}
我们还将诸凯勒形式组织成第二个组合，
\begin{equation} \label{eq:Omegazeta}
\hsymp(\zeta) = - \frac{i}{2\zeta} \omega_+ + \omega_3 - \frac{i}{2} \zeta \omega_-,
\end{equation}
其中我们引入了记号
\begin{equation}
\omega_\pm = \omega_1 \pm i \omega_2.
\end{equation}
$\hsymp(\zeta)$ 的本质性质是：对任一固定的 $\zeta \in \IC\IP^1$，它是 $\CM$ 上关于复结构 $J^\pz$ 的全纯辛形式。（要对 $\zeta = 0, \infty$ 理解这句话，须分别用 $\zeta$、$1/\zeta$ 去重标度 $\hsymp(\zeta)$。整体地，可以说 $\hsymp(\zeta)$ 被 $\IC\IP^1$ 上的线丛 $\CO(2)$ 扭转。）

\subsection{$g$ 的扭量构造}
\label{sec:twistor-general}

现在我们描述由 $\CM$ 上的全纯数据确定 $g$ 的方法；本文余下部分将一直使用它。

首先我们明确所作假设。回忆 $\CM$ 在拓扑上是 $\vm$ 上的环面丛。对 $\vm$ 中局部片、二次加细（quadratic refinement）、以及电荷格 $\clat$ 的局部截面 $\gamma$ 的每一种选取，我们都假定已给定一个局部定义的 $\IC^\times$ 值函数 $\CX_\gamma(u,\theta;\zeta)$，其自变量为 $(u,\theta) \in \CM$ 与 $\zeta \in \IC^\times$，并具有以下性质：
\begin{itemize}

\item 诸 $\CX_\gamma$ 是乘性的，
\begin{equation}
\CX_\gamma \CX_{\gamma'} = \CX_{\gamma + \gamma'}.
\end{equation}

\item 诸 $\CX_\gamma$ 满足一个实性条件，
\begin{equation} \label{eq:X-real}
\CX_\gamma(\zeta) = \overline{\CX_{-\gamma}(-1 / \bar\zeta)}.
\end{equation}

\item 所有 $\CX_\gamma$ 都是同一组微分方程的解，其形式为
\begin{align}
\frac{\partial}{\partial u^i} \CX &= \left( \frac{1}{\zeta} \CA^{(-1)}_{u^i} + \CA^{(0)}_{u^i} \right)  \CX , \label{eq:cr-1} \\
\frac{\partial}{\partial \bar{u}^{\bar i}} \CX &= \left( \CA^{(0)}_{\bar{u}^{\bar{i}}} + \zeta \CA^{(1)}_{\bar{u}^{\bar{i}}} \right) \CX, \label{eq:cr-2}
\end{align}
其中算子 $\CA^{(n)}_{u^i}$、$\CA^{(n)}_{\bar{u}^{\bar{i}}}$ 是环面纤维 $\CM_u$ 上的复竖直向量场，诸 $\CA^{(-1)}_{u^i}$ 在每一点线性无关，$\CA^{(1)}_{\bar{u}^{\bar{i}}}$ 亦然。（为说明这些方程的动机，注意我们在附录 B 中证明：$(\CM, g)$ 上的柯西--黎曼方程具有这种形式。）% TODO 合稿改 \ref{app:holomorphy}

\item 对每个固定的 $x \in \CM$，$\CX_\gamma(x; \zeta)$ 作为 $\zeta$ 的函数在 $\IC^\times$ 的一个稠密子集上全纯。（在下文的应用中，$\CX_\gamma(x; \zeta)$ 将在可数条直线之并以外全纯。）

\end{itemize}
为陈述关于诸函数 $\CX_\gamma$ 的最后三条假设，我们先定义
\begin{equation} \label{eq:darboux}
\hsymp(\zeta) := \frac{1}{8 \pi^2 R} \epsilon_{ij} \frac{d \CX_{\gamma^i}}{\CX_{\gamma^i}} \wedge \frac{d \CX_{\gamma^j}}{\CX_{\gamma^j}},
\end{equation}
其中 $d$ 指沿纤维的微分，即我们把 $\zeta$ 视为固定参数。我们假设：
\begin{itemize}

\item $\hsymp(\zeta)$ 是整体定义的（特别地，在 $\vm$ 的不同局部片上定义的 $\hsymp(\zeta)$ 彼此相符），并且关于 $\zeta \in \IC^\times$ 全纯。（注意这并不意味着诸 $\CX_\gamma$ 关于 $\zeta$ 全纯；在我们的应用中它们只是分段全纯的。）

\item $\hsymp(\zeta)$ 按全纯辛形式应有的意义非退化，即 $\ker \hsymp(\zeta)$ 是 $4r$ 维空间 $T_\IC \CM$ 中一个维数减半的子空间。

\item $\hsymp(\zeta)$ 当 $\zeta \to 0$ 或 $\zeta \to \infty$ 时只有单极点。

\end{itemize}
在本节余下部分，我们说明如何在 $\CM$ 上定义一个超凯勒度量 $g$，使得诸 $\CX_\gamma(\zeta)$ 是关于复结构 $J^\pz$ 的全纯函数，而 $\hsymp(\zeta)$ 就是 \eqref{eq:Omegazeta} 中那样的全纯辛形式。

我们考虑流形 $\CZ := \CM \times \IC\IP^1$。它具有以下性质：

\begin{enumerate}

\item  \emph{$\CZ$ 是复流形。}
在任一点 $(x, \zeta)$，$2r$ 个方程 \eqref{eq:cr-1}、\eqref{eq:cr-2} 定义了 $T_\IC \CM$ 的一个维数减半的子空间（若 $\zeta = 0$ 或 $\zeta = \infty$，把其中一个方程乘以因子 $\zeta$ 后此事仍成立）。该子空间与 $\partial / \partial \bar{\zeta}$ 所生成子空间的直和，是 $T_\IC \CZ$ 的一个维数减半的子空间。我们把这子空间定义为 $T^{0,1} \CZ$。这\emph{先验地}只定义了 $\CZ$ 上的一个近复结构。然而，诸函数 $\CX_\gamma$ 的存在保证了该近复结构实际上可积。（当然，诸 $\CX_\gamma$ 并非处处关于 $\zeta$ 全纯；但它们在稠密集上全纯，这已足以保证 Nijenhuis 张量消没。特别地，由此可知即使在 $\zeta = 0$ 或 $\zeta = \infty$ 附近，$\CZ$ 上也有复坐标。）

\item  \emph{$\CZ$ 是 $\IC\IP^1$ 上的全纯纤维化。}
投影就是 $p(x, \zeta) = \zeta$。

\item  \emph{存在 $\Omega^2_{\CZ/\IC \IP^1} \otimes \CO(2)$ 的一个全纯截面，它在每条纤维 $p^{-1}(\zeta)$ 上给出全纯辛形式。} 这就是整体定义的 $\hsymp(\zeta)$。

\item  \emph{存在一族全纯截面 $s: \IC\IP^1 \to \CZ$，每个的法丛都是 $N \simeq \CO(1)^{\oplus 2r}$。}
事实上，对每个 $x \in \CM$，可以定义截面 $s_x: \IC \IP^1 \to \CZ$ 为 $s_x(\zeta) = (x, \zeta)$。为看出它是全纯截面，先注意它至少在 $\zeta = 0,\infty$ 之外全纯，理由仅仅是 $\CZ$ 的局部复坐标 $\CX_{\gamma^i}(x, \zeta)$ 在固定 $x$ 时关于 $\zeta$ 全纯；但它连续延拓到 $\zeta = 0,\infty$，故由黎曼可去奇点定理，它在彼处也必全纯。为证明法丛 $N(s_x) \simeq \CO(1)^{\oplus 2r}$，先注意 $N^*(s_x)$ 的全纯截面与 $s_x$ 的一阶无穷小邻域上、在 $s_x$ 上取零的全纯函数之间有一一对应。但这样的函数由其在 $x$ 附近的一阶泰勒展开决定，即它们对应于平凡丛 $p^*((T^*_\IC)_x \CM)$ 的全纯截面，而这些截面零化由方程 \eqref{eq:cr-1}、\eqref{eq:cr-2} 定义的子丛 $B \subset p^*((T_\IC)_x \CM)$。取对偶，得 $N(s_x) \simeq p^*((T_\IC)_x \CM) / B$。另一方面，\eqref{eq:cr-1}、\eqref{eq:cr-2} 给出 $B \otimes \CO(1)$ 的 $2r$ 个平凡化截面。于是我们断定 $N(s_x) \otimes \CO(-1)$ 平凡。

\item  \emph{存在反全纯对合 $\sigma: \CZ \to \CZ$，它覆盖 $\IC \IP^1$ 上的对径映射，并按 $\sigma^* \hsymp = \overline{ \hsymp}$ 的意义保持 $\hsymp$。} 这个对合就是 $\sigma(x, \zeta) = (x, - 1 / \bar\zeta)$。利用实性条件 \eqref{eq:X-real}，可以验证它是反全纯的，且保持 $\hsymp$。
\end{enumerate}

这些正是一个超凯勒流形的\emph{扭量空间}的特征性质，如 \cite{Hitchin:1986ea,MR1206066} 所述。特别地，按 \cite{Hitchin:1986ea,MR1206066} 的配方，可以由 $\CZ$ 重构出 $\CM$ 上的超凯勒度量 $g$。$g$ 可以具体描述如下：注意由 $\hsymp(\zeta)^{r+1} = 0$ 可得 $\omega_+^r \wedge \omega_3 = 0$，后者蕴含实 2-形式 $\omega_3$ 在复结构 $J_3$ 下是 $(1,1)$ 型。因此我们可以用 $J_3$ 与 $\omega_3$ 在 $\CM$ 上造出一个凯勒度量 $g$。这个 $g$ 与扭量构造所保证的超凯勒度量一致。在随后几节中我们将采用这一进路。

\subsection{半平坦几何的扭量构造}
\label{sec:semiflat-twistor}

上述对超凯勒度量的描述，在半平坦度量 $g^\sfm$ 的情形尤为方便；该度量是我们在第 2.4 节引入的。% TODO 合稿改 \ref{sec:semiflat}
与前面一样，我们在 $\vm$ 的一个局部片上工作，并对二次加细作一个局部选取。于是对任一 $\gamma \in \clat$，我们写出局部定义的函数\footnote{这个公式最初是与鲍里斯·皮奥利纳（Boris Pioline）合作得到的；它本质上也是 \cite{Neitzke:2007ke} 中四元数凯勒情形的公式的刚性极限。它为发现本文所描述的构造提供了重要线索。}
\begin{equation} \label{eq:X-sf}
\CX_\gamma^\sfm(\zeta) := \exp\biggl[ \pi  R \zeta^{-1}  Z_\gamma + i
\theta_\gamma + \pi R \zeta \bar Z_\gamma \biggr].
\end{equation}
这些函数满足形如 \eqref{eq:cr-1}、\eqref{eq:cr-2} 的「柯西--黎曼方程」，其中
\begin{gather}
\CA^{(-1)}_{u^i} = -i \pi R \frac{\partial Z}{\partial u^i} \cdot \frac{\partial}{\partial \theta}, \quad \CA^{(1)}_{\bar{u}^{\bar{i}}} = -i \pi R \frac{\partial \bar{Z}}{\partial \bar{u}^{\bar{i}}} \cdot \frac{\partial}{\partial \theta}, \\
\CA^{(0)}_{u^i} = 0, \quad \CA^{(0)}_{\bar{u}^{\bar{i}}} = 0
\end{gather}
而 $Z$ 代表周期向量。于是 $\hsymp^\sfm(\zeta)$ 为
\begin{align} \label{eq:Omega-sf}
\hsymp^\sfm(\zeta) & := \frac{1}{8 \pi^2 R} \epsilon_{ij} \frac{d \CX^\sfm_{\gamma^i}}{\CX^\sfm_{\gamma^i}} \wedge \frac{d \CX^\sfm_{\gamma^j}}{\CX^\sfm_{\gamma^j}} \\
&= \frac{1}{4\pi} \left[ \frac{i}{\zeta} \left\langle d Z, d \theta \right\rangle + \left( \pi R \langle d Z, d \bar{Z} \rangle - \frac{1}{2 \pi R} \langle d \theta , d \theta \rangle \right) + i \zeta \langle d \bar{Z}, d \theta \rangle \right]. \label{eq:Omega-sf-2}
\end{align}
（注意：消没条件式 (2.5) 保证 $\hsymp^\sfm(\zeta)$ 不含 $\zeta^{-2}$ 阶与 $\zeta^2$ 阶的项。）% TODO 合稿改 \eqref{eq:lag-constraint}
$\hsymp^\sfm(\zeta)$ 与 $\CX^\sfm_\gamma(\zeta)$ 满足我们在第 \ref{sec:twistor-general} 节所述构造所需的诸条件，因此它们就是 $\CM$ 上某个超凯勒度量的全纯辛形式与复坐标。正如我们下面要验证的，这个度量正是所求的 $g^\sfm$。

先注意，比较 \eqref{eq:Omegazeta} 与 \eqref{eq:Omega-sf-2} 中的领头项，给出
\begin{equation} \label{eq:omegaplus-sf}
\omega_+^\sfm = - \frac{1}{2\pi} \inprod{d Z, d \theta}.
\end{equation}
由 $\omega_+^\sfm$ 可以确定复结构 $J^\sfm_3$：事实上，选定一个电磁对偶标架后，我们可以把 \eqref{eq:omegaplus-sf} 改写为
\begin{equation}
\omega_+^\sfm = \frac{1}{2\pi} da^I \wedge dz_I.
\end{equation}
这就清楚地表明，$\CM$ 在复结构 $J^\sfm_3$ 之下正是由复环面构成的 \SW 纤维化。这正是我们在第 2.4 节已经描述过的复结构。% TODO 合稿改 \ref{sec:semiflat}

类似地，比较 \eqref{eq:Omegazeta} 与 \eqref{eq:Omega-sf-2} 中与 $\zeta$ 无关的项，给出
\begin{equation} \label{eq:omega3-sf}
\omega_3^\sfm = \frac{R}{4} \langle d Z , d \bar{Z} \rangle - \frac{1}{8 \pi^2 R} \langle d\theta, d \theta \rangle,
\end{equation}
它可改写为
\begin{equation}
\omega_3^\sfm = \frac{i}{2} \left( R (\im \tau)_{IJ} d a^I \wedge d \bar{a}^J + \frac{1}{4 \pi^2 R} ((\im \tau)^{-1})^{IJ} dz_I \wedge d\bar{z}_J \right).
\end{equation}
将此式与式 (2.35) 比较可知，$g^\sfm$ 确实关于复结构 $J^\sfm_3$ 是凯勒的、且凯勒形式为 $\omega_3^\sfm$，因而是由 $\hsymp^\sfm(\zeta)$ 出发的扭量构造所保证的那个超凯勒度量。% TODO 合稿改 \eqref{eq:semiflatmetric}

本节我们看到，$\CM$ 上的半平坦度量及其超凯勒结构，都可以由 \eqref{eq:X-sf} 中定义的诸函数 $\CX_\gamma^\sfm$ 构造出来。这些函数对下文具有根本的重要性。

\begin{yizhu}
$\hsymp(\zeta)$ 在原文中出现两次定义：式~\eqref{eq:Omegazeta} 用它指超凯勒结构给出的全纯辛形式，式~\eqref{eq:darboux} 又以「$:=$」把它定义为由诸 $\CX_\gamma$ 组装出的量。两者是同一个对象：本节的构造正是要求由式~\eqref{eq:darboux} 得到的 $\hsymp(\zeta)$ 恰好等于式~\eqref{eq:Omegazeta} 中的那一个（见本节末的陈述）。
\end{yizhu}
